This section establishes the rank-$r$ rate-distortion lower
bound. The argument rests on three ingredients: a deterministic
max-$\geq$-avg identity that converts max-distortion into an
average-distortion problem plus an irreducible floor
(\S\ref{subsec:id}); a closed-form evaluation of that floor under
the hard distribution $P^\star$
(\S\ref{subsec:floor}); and a rate-distortion lower bound on the
remaining average-distortion subproblem via a Shannon converse
applied conditionally on the curvature data $\boldsymbol H$
(\S\ref{subsec:main}).

\subsection{The $H_t$-weighted max $\geq$ avg identity}
\label{subsec:id}

The Phase~0 identity $T^{-1}\sum_t\|w-\tau_t\|^2 = \|w-\bar\tau
\|^2 + T^{-1}\sum_t\|\tau_t-\bar\tau\|^2$ relies on
$\sum_t(\tau_t-\bar\tau) = 0$, which depends only on $\bar\tau$
being the arithmetic mean. The right generalization to
$H_t$-weighted distortion uses $\tauH$ of~\eqref{eq:derived} as
the centroid.

\begin{lemma}[$H_t$-weighted max $\geq$ avg identity]\label{lem:id}
For every admissible $(\boldsymbol\tau, \boldsymbol H)$ and every
$w\in\R^d$,
\begin{equation}\label{eq:id-eq}
  \tfrac1T\sum_{t=1}^T (w-\tau_t)^\top H_t (w-\tau_t)
  \;=\; (w-\tauH)^\top\bar H(w-\tauH) + \tfrac1T\sum_{t=1}^T
  (\tau_t-\tauH)^\top H_t (\tau_t-\tauH),
\end{equation}
and hence
$\max_t (w-\tau_t)^\top H_t(w-\tau_t) \geq
(w-\tauH)^\top\bar H(w-\tauH) + T^{-1}\sum_t (\tau_t-\tauH)^\top
H_t (\tau_t-\tauH)$.
\end{lemma}

The full proof is given in the supplementary material. The
argument is a two-line algebraic expansion: the cross-term
$\sum_t H_t(\tau_t-\tauH) = T(\bar H\tauH - \bar H\tauH) = 0$
vanishes by the defining property $\bar H\tauH = T^{-1}\sum_t
H_t\tau_t$ of the centroid. No probabilistic argument is needed;
at $H_t = I_d$ we recover the Phase~0 identity exactly.

Two consequences of Lemma~\ref{lem:id} drive everything that
follows. First, \emph{components of $w$ orthogonal to
$\range(\bar H) = V_1+\cdots+V_T$ annihilate against every $H_t$
and do not affect~\eqref{eq:id-eq}}; without loss of generality
the decoder's output lies in an at-most-$\deff$-dimensional
subspace. Second, the max-distortion lower bound separates into a
compression term in $\bar H$-metric on $\tauH$ and an irreducible
floor term that does not depend on $w$ at all. The rest of the
argument bounds each piece.

\subsection{Closed-form floor}\label{subsec:floor}

Under the hard distribution $P^\star$ of \S\ref{sec:setup}, the
floor term admits an exact closed form depending only on
$\deff$.

\begin{lemma}[Closed-form floor]\label{lem:floor-general}
Under $P^\star$ with any fixed positive-definite spectra
$(D_t)_{t=1}^T$ (possibly distinct across tasks),
\begin{equation}\label{eq:floor-closed}
  \E_{P^\star}\!\left[\tfrac1T\sum_{t=1}^T(\tau_t-\tauH)^\top
   H_t(\tau_t-\tauH)\right]
  \;=\; B^2\!\left(1 - \frac{\E[\deff]}{Tr}\right).
\end{equation}
\end{lemma}

The proof (supplementary material) splits into a deterministic
identity reducing the floor to a bilinear form in $\tauH$ and an
expectation-via-isotropy step. The key probabilistic input is
the second-moment identity $\E[H_t\tau_t\tau_t^\top H_t\mid U_t]
= (B^2/r) H_t$, which holds \emph{for any} $D_t$ because of the
$D_t^{-1/2}$ stretch in the definition of $P^\star$.

Three limiting cases illuminate the formula.

\begin{enumerate}[leftmargin=*,topsep=2pt,itemsep=2pt]
  \item \textbf{Full-rank limit} ($r = d$, $V_t = \R^d$). Then
  $\deff = d$ a.s., floor $= B^2(1 - d/(Td)) = B^2(1 - 1/T)$,
  recovering the Phase~0 formula for classical averaged-vector
  quantization.
  \item \textbf{Single-task limit} ($T = 1$). Then $\deff = r$,
  floor $= B^2(1 - r/r) = 0$, and the rate-distortion function
  reduces to single-source quantization of $\tau_1\in V_1\cong
  \R^r$. This case matches TurboQuant's Theorem~3
  \cite{zandieh2025turboquant}.
  \item \textbf{Shared-subspace limit} ($V_t = V$ for all $t$).
  Then $\deff = r$, floor $= B^2(1 - r/(Tr)) = B^2(1-1/T)$, again
  recovering Phase~0 with the ambient dimension replaced by the
  rank.
\end{enumerate}
In the generic LoRA regime ($V_t$ Stiefel-random and $Tr\leq d$),
$\deff = Tr$ almost surely and the floor vanishes. The
rate-distortion function in this regime is controlled entirely by
the compression term.

\subsection{Main lower bound}\label{subsec:main}

Combining Lemma~\ref{lem:id} and Lemma~\ref{lem:floor-general}
with a rank-$r$ analog of the Phase~0 rotational-invariance
argument gives the main theorem.

\begin{theorem}[Rank-$r$ RD lower bound, general $H_t$]
\label{thm:general}
Under $P^\star$ with any fixed positive-definite spectra
$(D_t)_{t=1}^T$, and with $\rank(H_t) = r$ for all $t$,
\begin{equation}\label{eq:thm-lb}
  \mathcal{D}^\star(T, d, r, B, R)
  \;\geq\;
  B^2\!\left(1 - \frac{\E[\deff]}{Tr}\right)
  \;+\; c_{\mathrm{TQ}}\cdot\frac{B^2\E[\deff]}{Tr}
  \cdot 2^{-2R/\E[\deff]}.
\end{equation}
Under Stiefel-random sampling with $Tr\leq d$, $\deff = Tr$ a.s.,
the floor vanishes, and
\begin{equation}\label{eq:thm-lb-stiefel}
  \mathcal{D}^\star(T,d,r,B,R)
  \;\geq\;
  c_{\mathrm{TQ}}\cdot B^2 \cdot 2^{-2R/(Tr)},
\end{equation}
with $c_{\mathrm{TQ}}$ the universal constant from
TurboQuant~\cite{zandieh2025turboquant} Theorem~3.
\end{theorem}

\paragraph{Proof sketch.} The argument is a six-step reduction.

\emph{Step 1 (Yao's minimax).} $\mathcal{D}^\star\geq
\inf_{E,D}\E_{P^\star}[\max_t D_t(w^\star)]$ by
\eqref{eq:yao}; all subsequent steps bound the RHS.

\emph{Step 2 (max $\geq$ avg reduction).} By Lemma~\ref{lem:id}
and Lemma~\ref{lem:floor-general},
\begin{equation*}
  \E_{P^\star}[\max_t D_t(w^\star)]
  \;\geq\;
  \E_{P^\star}\bigl[(w^\star-\tauH)^\top\bar H(w^\star-\tauH)\bigr]
  + B^2(1 - \E[\deff]/(Tr)).
\end{equation*}

\emph{Step 3 (restrict to $\range(\bar H)$).} By the orthogonality
remark following Lemma~\ref{lem:id}, the decoder may be assumed
to emit $w^\star\in\range(\bar H)$ without loss.

\emph{Step 4 (change of variable).} Let $\eta := \bar H^{1/2}
\tauH$, viewed as an element of $\range(\bar H)\cong
\R^{\deff}$. Then $(w^\star-\tauH)^\top\bar H(w^\star-\tauH) =
\|\bar H^{1/2}w^\star - \eta\|_{\R^{\deff}}^2$, so the
compression subproblem is rate-$R$ quantization of $\eta$ under
squared Euclidean loss in $\R^{\deff}$.

\emph{Step 5 (Shannon converse, conditional on $\boldsymbol H$).}
Because both ends may use $\boldsymbol H$ as side information
(\S\ref{sec:setup}), we bound the compression term conditional on
the full curvature data $\boldsymbol H$ rather than on the
subspace $W := \range(\bar H)$ alone. Given $\boldsymbol H$, the
only remaining randomness is in the uniform directions $u_t$, so
$\eta = \bar H^{1/2}\tauH$ is a zero-mean source on
$\range(\bar H)\cong\R^{\deff}$. The Shannon rate-distortion
converse for a $\deff$-dimensional source under squared-error loss
holds for \emph{every} rate-$R$ code, including decoders and
codebooks adapted to $\boldsymbol H$, and gives
\begin{equation*}
  \E\bigl[\|\bar H^{1/2}w^\star-\eta\|^2\,\bigm|\,\boldsymbol H\bigr]
  \;\geq\;
  c_{\mathrm{TQ}}\cdot\E\bigl[\|\eta\|^2\,\bigm|\,\boldsymbol H\bigr]\cdot
  2^{-2R/\deff}.
\end{equation*}
Conditioning on $\boldsymbol H$ — rather than averaging over
$O(W)$ rotations — is what makes the converse valid against an
$\boldsymbol H$-aware decoder such as the achievability
construction of \S\ref{sec:achv}. The constant $c_{\mathrm{TQ}}$
now also absorbs the entropy-power-to-variance ratio of $\eta$ and
the condition number of its conditional covariance, both $O(1)$
under $P^\star$ because the spectra $D_t$ are fixed and
$\tau_t^\top H_t\tau_t = B^2$ fixes the scale. The floor's
$D_t$-freeness (Lemma~\ref{lem:floor-general}) is unaffected: for
each fixed spectrum $v_t = B D_t^{-1/2} u_t$ is rotation-uniform
in $\R^r$, so the constant is uniform over the admissible $D_t$.

\emph{Step 6 (combine).} Taking expectation over $\boldsymbol H$ and using
$\E[\|\eta\|^2] = B^2\E[\deff]/(Tr)$ (from
Lemma~\ref{lem:floor-general}), substitution into Step~2 yields
\eqref{eq:thm-lb}. The Stiefel specialization $Tr\leq d$ gives
$\deff = Tr$ a.s., hence \eqref{eq:thm-lb-stiefel}.

The full proof is given in the supplementary material; the
six-step skeleton above is the complete argument modulo the
algebraic expansions in Steps~2 and~6 and the Shannon-bound
machinery of Step~5.

\paragraph{Why the general-$H_t$ statement is the headline.}
Real LoRA fine-tunes have task-specific Hessians with
non-trivial spectra. The Fisher information at a fine-tuned
point $\tau_t$ is rank-$r$ PSD with range in $V_t$ but generally
not a projector; its spectrum inside $V_t$ is the task-$t$
curvature and generally differs across tasks.
Theorem~\ref{thm:general} applies directly to this setting
without any modeling assumption on the $D_t$'s beyond
positive-definiteness. This closes the $\mathrm{MSE}\to
\mathrm{CE}$ bridge for the lower-bound half of the
rate-distortion characterization; the achievability analysis in
\S\ref{sec:achv} proceeds under the same general hypothesis.
